\clearpage
\section{补充训练与执行对照}
\label{app:additional_zh}

\subsection{双seed训练配方比较}
该实验比较开发候选G（冻结DINOv3-S）与I（可训练CNN）；二者均保留各自P07监督、完整视觉窗口和当前目标RGB重建。每个配置使用训练seed0和1，比较十六卡global batch16与三十二卡global batch64。大batch同时将学习率从$2\times10^{-4}$提高至$4\times10^{-4}$，warm-up从1,000改为250更新，仍对应16,000个样本。因此结果反映配方和资源的共同变化，不能单独归因于batch。

八次训练均完成四轮、641,664个样本。B16执行40,104次更新，B64执行10,028次；每轮32个尾部样本保留，不填充、重复或丢弃。AdamW、梯度裁剪、FP32/TF32、损失、归一化和按四轮样本进度设置的cosine调度保持一致。所有最终端点统一采用十六卡full350验证和原1,400场评测，场景及推理seed配对。B16的seed0为从头重跑，逐场成败与相应BASE2端点完全一致，但不增加独立训练seed数；seed1才是新增重复。

\begingroup
\renewenvironment{table*}[1][]{\begin{readingtable}}{\end{readingtable}}
% Generated from Research BATCH0/results/full_measurements_20261016.json.
\begin{table*}[t]
\centering\small
\caption{\tablelang{Completed training-recipe study, four epochs per run. G uses frozen
DINOv3-S and I a trainable CNN; both use the corresponding P07 combination.
Batch 64 jointly changes learning rate and device count (details in text).
Training time excludes saving and validation. Job time runs from job creation
through complete evaluation, including the observed queue and startup.}{已完成的训练配方研究，每次训练四轮。G 使用冻结的 DINOv3-S，I 使用可训练 CNN；二者均采用对应的 P07 组合。Batch 64 同时改变学习率和设备数（详见正文）。训练耗时不含保存和验证。作业耗时从创建作业起计至完整评测结束，包含实际发生的排队和启动时间。}}
\label{tab:batch_recipe}
\begin{tabular}{@{}lrrrrrrr@{}}
\toprule
\tablelang{Profile}{配置} & \tablelang{Seed}{种子} & \tablelang{Batch}{批大小} & \tablelang{GPUs}{GPU 数} & \tablelang{Success /1,400}{成功数 /1,400} & \tablelang{SR (\%)}{成功率 (\%)} & \tablelang{Train (min)}{训练 (min)} & \tablelang{Job (min)}{作业 (min)} \\
\midrule
G & 0 & 16 & 16 & 849 & 60.64 & 263.40 & 301.28 \\
G & 0 & 64 & 32 & 855 & 61.07 & 99.27 & 121.32 \\
G & 1 & 16 & 16 & 834 & 59.57 & 263.28 & 303.58 \\
G & 1 & 64 & 32 & 822 & 58.71 & 97.48 & 120.23 \\
I & 0 & 16 & 16 & 935 & 66.79 & 220.75 & 257.20 \\
I & 0 & 64 & 32 & 884 & 63.14 & 78.86 & 101.74 \\
I & 1 & 16 & 16 & 857 & 61.21 & 224.91 & 258.42 \\
I & 1 & 64 & 32 & 935 & 66.79 & 79.80 & 101.02 \\
\bottomrule
\end{tabular}
\end{table*}

% Generated from the 11,200 original BATCH0 episode results.
\begin{table*}[t]
\centering\small
\caption{\tablelang{Training-recipe successes per task (100 scenes per entry).
B16 and B64 include their respective learning-rate/device settings;
s0/s1 denote training seeds. All runs use the same accepted scene identities
and inference seeds. Differences across tasks remain visible despite similar totals.}{训练配方研究的逐任务成功数（每项 100 个场景）。B16 和 B64 分别包含各自的学习率与设备设置；s0/s1 表示训练 seed。所有运行使用相同的已接受场景身份和推理 seed。总分相近时，逐任务差异仍然存在。}}
\label{tab:batch_tasks}
\begin{tabular}{@{}lrrrrrrrr@{}}
\toprule
 & \multicolumn{4}{c}{\tablelang{G: frozen DINOv3-S}{G：冻结 DINOv3-S}} & \multicolumn{4}{c}{\tablelang{I: trainable CNN}{I：可训练 CNN}} \\
\cmidrule(lr){2-5}\cmidrule(l){6-9}
\tablelang{Task}{任务} & B16/s0 & B64/s0 & B16/s1 & B64/s1 & B16/s0 & B64/s0 & B16/s1 & B64/s1 \\
\midrule
\tablelang{Grasp}{抓取} & 92 & 87 & 91 & 89 & 90 & 95 & 87 & 96 \\
\tablelang{Belt}{传送带} & 100 & 100 & 100 & 97 & 97 & 100 & 99 & 100 \\
\tablelang{Occlusion}{遮挡} & 100 & 99 & 100 & 100 & 97 & 98 & 97 & 100 \\
SO(3) & 87 & 88 & 68 & 53 & 86 & 75 & 56 & 55 \\
XYZ & 75 & 58 & 60 & 69 & 82 & 69 & 77 & 85 \\
\tablelang{SE(3) placement}{SE(3) 放置} & 0 & 0 & 0 & 0 & 0 & 0 & 1 & 0 \\
\tablelang{Trans. placement}{平移放置} & 1 & 10 & 10 & 4 & 4 & 5 & 13 & 29 \\
\tablelang{Full passage}{完整穿越} & 0 & 0 & 0 & 0 & 0 & 0 & 0 & 0 \\
\tablelang{Phase passage}{相位穿越} & 0 & 0 & 13 & 0 & 72 & 32 & 14 & 46 \\
\tablelang{Can}{罐体} & 99 & 99 & 97 & 98 & 100 & 98 & 99 & 100 \\
\tablelang{Slider}{滑块} & 83 & 81 & 76 & 80 & 81 & 86 & 75 & 89 \\
\tablelang{Sphere}{球体} & 40 & 50 & 50 & 57 & 56 & 47 & 60 & 73 \\
\tablelang{Tabletop sphere}{桌面球体} & 100 & 99 & 99 & 100 & 98 & 99 & 99 & 100 \\
\tablelang{Replan}{重规划} & 72 & 84 & 70 & 75 & 72 & 80 & 80 & 62 \\
\midrule
\tablelang{Total}{总计} & 849 & 855 & 834 & 822 & 935 & 884 & 857 & 935 \\
\bottomrule
\end{tabular}
\end{table*}

\endgroup

G在seed0/1下，B64相对B16分别变化$+0.43$和$-0.86$个百分点，配对95\%区间为$[-1.21,2.07]$、$[-2.50,0.79]$。I的变化方向相反：$-3.64$ $[-5.64,-1.64]$和$+5.57$ $[3.57,7.57]$个百分点。单看B16，I从seed0的935场成功降至seed1的857场，相位穿越从72场降至14场。原seed0高分和配方平均收益均不能证明稳定性；匹配控制E/F尚未完成seed1。

两seed的平均配方差值为G的$-0.21$ $[-1.36,0.96]$和I的$+0.96$ $[-0.50,2.39]$个百分点。区间使用20,000次百分位bootstrap、随机种子20261017；每任务从100个场景身份中重采样，同时保留batch与训练seed配对。区间以这两个已训练seed和十四任务为条件，不估计新训练seed的变化。逐任务表保留总体分数掩盖的损失：G的seed0在XYZ少成功17场，seed1在SO(3)少15场。辅助指标也未全面保持，例如I的位姿旋转误差分别增加6.17\%和14.35\%。

从任务创建到完整结果的平均耗时，G由302.43降至120.77分钟，I由257.81降至101.38分钟。这些实测包含排队、准备、训练、保存、验证、导出和评测，不是单独的硬件加速率。训练计算GPU小时分别由70.22降至52.47、59.42降至42.31；该口径排除保存和验证，也不是平台账单。预先设定的总体SR工作容忍界为下降3个百分点，G的两个已测seed满足，I不满足。因此通用配方保留B16；B64只作为当前G、以总体SR为准时的加速选项，不推广到I、单任务敏感比较或尚未验证的最终设计。

\subsection{冻结策略的三态夹爪读出}
\label{app:learned_decoder_zh}
该实验冻结第五轮P03策略，学习开／保持／关读出。输入共889维，拼接最后一次Euler评估（$\tau=0.8$）的当前动作隐状态、八帧已有EEF位姿／夹爪开度／有效掩码、采样出的六维机械臂动作序列和当前原始夹爪值。隐层宽度为128的GELU MLP输出三类logits，argmax只替换实际执行的slot0夹爪命令。raw-only对照使用相同隐层宽度，只读取当前原始夹爪值。两者不增加接触真值或未来标签输入；原策略、归一化、五步Euler与同步单步执行保留。

原11,200条训练和350条验证示范分别有1,324,374和41,533个有效下发命令标签；每个任务内，训练与验证的示范source seed无交集。每个有效当前命令都进入样本，包括后续动作槽已无效的尾部窗口。保持标签只出现在接球：训练2,562个，占全部标签0.1934\%；验证79个。其余十三任务没有有效保持标签，相位穿越的训练和验证标签全部为关。因此，这些标签没有覆盖其他任务可能需要的保持动作。

两头均用seed0训练二十轮，AdamW学习率$3\times10^{-4}$、weight decay $10^{-4}$、batch4,096、梯度裁剪1，普通交叉熵且无类别重加权。特征均值和方差只用训练集。按验证集三类NLL的平均值选最低端点，同值取最早；context选中第15轮、NLL为0.01330，raw-only选中第20轮、NLL为0.19640。

\begingroup
\renewenvironment{table}[1][]{\begin{readingtable}}{\end{readingtable}}
% Generated from Research GRIPPER2/results/learned/training-report.json.
\begin{table}[t]
\centering\small
\caption{Frozen-decoder validation confusion matrices: 41,533 valid issued
commands from 350 demonstrations. Rows are ground truth; columns are predicted
open/hold/close. These are offline classification counts, not policy successes.}
\label{tab:decoder_confusion}
\begin{tabular}{@{}lrrrrrr@{}}
\toprule
 & \multicolumn{3}{c}{Raw only} & \multicolumn{3}{c}{Context} \\
\cmidrule(lr){2-4}\cmidrule(l){5-7}
GT & Open & Hold & Close & Open & Hold & Close \\
\midrule
Open & 8003 & 1 & 214 & 8132 & 0 & 86 \\
Hold & 0 & 75 & 4 & 0 & 79 & 0 \\
Close & 135 & 9 & 33092 & 105 & 4 & 33127 \\
\bottomrule
\end{tabular}
\end{table}

\endgroup

context的保持recall为79/79、precision为79/83，raw-only为75/79和75/85。153个一至两tick短脉冲中，两者分别正确预测137和93个，均没有误判为保持。这些结果来自示范观测上的冻结特征，不能当作闭环策略成绩。
三类NLL均有改善。79个保持标签在端点选择指标中占三分之一，贡献该指标降幅的74.5\%；训练交叉熵仍未加权。context在验证集仍有86次开变闭、105次闭变开错误，因此平均损失不约束单次闭环决策的代价。

预先确定的十一例开发集回归案例中，context通过E5接球0/9/14/80/85、E5重规划49/88，以及E7 XYZ64与重规划7；E7直接使用E5训练的头，不再拟合。失败为E5接球52和重规划35。9/11未达到全部十一例成功的条件，因此未执行原计划140个新场景乘三种解码器的420场评测，也未采用该读出。开发集回归清单的通过数不是独立成功率。

两个公开动作API在观察模式下保持原命令及RNG；二十八场sign／0.1预检复现原结果。本地重算608条下发命令，CPU logits最大差为$3.82\times10^{-6}$，argmax命令完全一致；其余六维动作与预测序列也通过记录对照。完整逐帧策略输入张量未保存，不能据此声称输入逐项完全相同。

图~\ref{fig:decoder_failures_zh}展示全部两例失败。两者首次干预均把约$-1$的开爪翻成$+1$闭爪，不只是删除近零输出。接球在0.650\,s改变命令，0.666\,s开度分叉，之后没有双侧接触；重规划经过原生一tick队列，于1.550\,s应用新命令，1.556\,s开度分叉。候选分别于1.324\,s接触托盘、5.174\,s越过下游边界而失败；sign两例均成功，固定0.1在重规划成功、接球同样失败。干预前已记录的预测、命令与2\,ms物理字段和对照一致，后续机械臂与物体分叉包含闭环反馈。这些记录解释该读出的失败，不能代替未来预测到动作接口的干预实验。

\begin{figure}[ht]
\centering
\includegraphics[width=\textwidth]{decoder_failures.pdf}
\caption{冻结夹爪读出的两例失败，来自预先确定的开发集回归案例。命令与实际开度使用物理时间，事件条保留各运行的接触和终止结果。sign、固定0.1死区与context头共用冻结策略；头的首次干预将开爪翻为闭爪，随后出现物理分叉。局部曲线与全程事件采用各自标注的时间范围，raw预测为20\,Hz，物理采样间隔2\,ms；分叉指learned相对sign的开度或逐坐标位置差超过10\,$\mu$m。双侧碰撞接触与重规划附着分开表示。该图不提供新场景成功率或未来预测干预结论。}
\label{fig:decoder_failures_zh}
\end{figure}

\paragraph{已保存特征的敏感性。}
在接球52和重规划35的首次翻转处，仅换用前一帧动作隐状态，闭类减开类logit由$+2.452/+0.700$变为$-13.839/-9.811$；将当前隐状态换入前一帧其余特征，则得到$+2.243/+3.184$。仅替换原始夹爪值不能消除翻转。这些混合特征检验固定分类器的局部响应，不一定对应可实现的观测或成功执行。首次翻转发生在该头此前尚未干预时；保留数据仍不能区分示范到策略的特征失配与末次Euler隐状态的采样敏感性，两种解释均待验证。回归未通过的结论不变。
